\section*{Irrationality of $\sqrt{7}$}

\noindent\textbf{Definition.}
A real number $r$ is \emph{irrational} if there do not exist integers $p, q$ with $q \neq 0$ such that $r = p/q$.

\bigskip

We first record three elementary lemmas that will be used in the proof.

\begin{lemma}[Divisibility of a difference]\label{lem:div-sub}
If $c, x, y \in \mathbb{Z}$, $c \mid x$, and $c \mid y$, then $c \mid (x - y)$.
\end{lemma}

\begin{proof}
Since $c \mid x$ and $c \mid y$, there exist integers $u, v$ such that $x = cu$ and $y = cv$.
Therefore $x - y = cu - cv = c(u - v)$. Since $u - v \in \mathbb{Z}$, we have $c \mid (x - y)$.
\end{proof}

\begin{lemma}[Finite maximum]\label{lem:fin-max}
If $N \in \mathbb{Z}_{>0}$ and $S$ is a nonempty subset of $\{1, 2, \ldots, N\}$, then $S$ has a largest element.
\end{lemma}

\begin{proof}
We prove this by induction on $N$.

\textit{Base case.} For $N = 1$, the only nonempty subset of $\{1\}$ is $\{1\}$, which has largest element $1$.

\textit{Inductive step.} Assume the result holds for some $N \geq 1$. Let $S$ be a nonempty subset of $\{1, 2, \ldots, N, N+1\}$. If $N+1 \in S$, then $N+1$ is the largest element of $S$. If $N+1 \notin S$, then $S \subseteq \{1, \ldots, N\}$, so by the induction hypothesis $S$ has a largest element.

By induction, the lemma holds for all $N \in \mathbb{Z}_{>0}$.
\end{proof}

\begin{lemma}[Lowest-terms reduction]\label{lem:low-terms}
If $p, q \in \mathbb{Z}$ and $q \neq 0$, then there exist integers $a, b$ with $b \neq 0$ such that $p/q = a/b$ and $a, b$ have no positive common divisor greater than $1$.
\end{lemma}

\begin{proof}
Define $D = \{d \in \mathbb{Z}_{>0} : d \mid p \text{ and } d \mid q\}$. Since $1 \mid p$ and $1 \mid q$, we have $1 \in D$, so $D$ is nonempty.

For any $d \in D$, since $d \mid q$ there exists $t \in \mathbb{Z}$ with $q = dt$. Since $q \neq 0$, we have $t \neq 0$, so $|t| \geq 1$. Thus $|q| = d|t| \geq d$, giving $d \leq |q|$. Hence $D \subseteq \{1, 2, \ldots, |q|\}$. By Lemma~\ref{lem:fin-max}, $D$ has a largest element; call it $d$.

Since $d \mid p$ and $d \mid q$, write $p = da$ and $q = db$ for integers $a, b$. Then $b \neq 0$ (since $q \neq 0$) and $p/q = a/b$.

We claim $a$ and $b$ have no positive common divisor greater than $1$. Suppose for contradiction that $c > 1$ with $c \mid a$ and $c \mid b$. Write $a = ca_1$ and $b = cb_1$. Then $p = dca_1$ and $q = dcb_1$, so $dc \mid p$ and $dc \mid q$. Since $d > 0$ and $c > 1$, we have $dc > d$ and $dc \in D$, contradicting the maximality of $d$.
\end{proof}

\begin{lemma}[Prime divisibility for $7$]\label{lem:prime7}
For any $n \in \mathbb{Z}$, if $7 \mid n^2$ then $7 \mid n$.
\end{lemma}

\begin{proof}
By the division algorithm, write $n = 7a + r$ with $a \in \mathbb{Z}$ and $r \in \{0, 1, 2, 3, 4, 5, 6\}$. Then
\[
n^2 = (7a + r)^2 = 49a^2 + 14ar + r^2,
\]
so $n^2 - r^2 = 7(7a^2 + 2ar)$, hence $7 \mid (n^2 - r^2)$.

By hypothesis $7 \mid n^2$, so Lemma~\ref{lem:div-sub} gives $7 \mid (n^2 - (n^2 - r^2)) = r^2$.

Checking all residues: $r^2 \in \{0, 1, 4, 9, 16, 25, 36\}$ for $r \in \{0,\ldots,6\}$. Of these, only $r^2 = 0$ is divisible by $7$, so $r = 0$. Therefore $n = 7a$, giving $7 \mid n$.
\end{proof}

\begin{theorem}\label{thm:sqrt7-irr}
$\sqrt{7}$ is irrational.
\end{theorem}

\begin{proof}
Suppose for contradiction that $\sqrt{7}$ is not irrational. Then by definition there exist integers $p, q$ with $q \neq 0$ such that $\sqrt{7} = p/q$.

By Lemma~\ref{lem:low-terms}, we may write $\sqrt{7} = a/b$ where $a, b \in \mathbb{Z}$, $b \neq 0$, and $a$ and $b$ have no positive common divisor greater than $1$.

Squaring gives $7 = a^2/b^2$, and multiplying both sides by $b^2$ (valid since $b \neq 0$):
\[
7b^2 = a^2.
\]
Hence $7 \mid a^2$. By Lemma~\ref{lem:prime7}, $7 \mid a$, so $a = 7k$ for some $k \in \mathbb{Z}$.

Substituting: $7b^2 = (7k)^2 = 49k^2$. Dividing by $7$: $b^2 = 7k^2$, hence $7 \mid b^2$. By Lemma~\ref{lem:prime7} again, $7 \mid b$.

Thus $7 \mid a$ and $7 \mid b$. Since $7 > 1$, the integer $7$ is a positive common divisor of $a$ and $b$ greater than $1$, contradicting our choice of $a$ and $b$.

Therefore $\sqrt{7}$ is irrational.
\end{proof}
